\documentclass[12pt]{article}

\usepackage[T1]{fontenc}
\usepackage[utf8]{inputenc}
\usepackage[romanian]{babel}

\usepackage{amsmath, amssymb, amsthm}
\usepackage{geometry}
\geometry{a4paper, margin=2.5cm}

\title{Media, Varianta și Covarianța Variabilelor Aleatoare}
\author{}
\date{}

\theoremstyle{plain}
\newtheorem{definition}{Definiție}
\newtheorem{proposition}{Propoziție}
\newtheorem{example}{Exemplu}

\begin{document}

\maketitle

\section{Media}

\subsection{Caz discret}

\begin{definition}
Fie $X$ o variabilă aleatoare discretă cu valori $x_k$ și funcție de masă $p(x_k)$. Atunci


\[
\mathbb{E}[X] = \sum_k x_k\, p(x_k),
\]


dacă seria $\sum_k |x_k| p(x_k)$ este convergentă.
\end{definition}

\subsection{Caz continuu}

\begin{definition}
Fie $X$ o variabilă aleatoare continuă cu densitate $f$. Atunci


\[
\mathbb{E}[X] = \int_{-\infty}^{\infty} x f(x)\, dx,
\]


dacă integrala $\int |x| f(x)\, dx$ este finită.
\end{definition}

\section{Varianta}

\subsection{Definiția generală}

\begin{definition}
Pentru o variabilă aleatoare $X$ cu media $\mu = \mathbb{E}[X]$, varianta este


\[
\mathrm{Var}(X) = \mathbb{E}[(X - \mu)^2].
\]


\end{definition}

\begin{proposition}[Formula dezvoltată]


\[
\mathrm{Var}(X) = \mathbb{E}[X^2] - (\mathbb{E}[X])^2.
\]


\end{proposition}

\subsection{Caz discret}

Dacă $X$ ia valori $x_k$ cu probabilități $p(x_k)$, atunci


\[
\mathrm{Var}(X)
= \sum_k (x_k - \mu)^2 p(x_k)
= \sum_k x_k^2 p(x_k) - \left(\sum_k x_k p(x_k)\right)^2.
\]



\subsection{Caz continuu}

Dacă $X$ are densitate $f$, atunci


\[
\mathrm{Var}(X)
= \int_{-\infty}^{\infty} (x - \mu)^2 f(x)\, dx
= \int x^2 f(x)\, dx - \left(\int x f(x)\, dx\right)^2.
\]



\section{Covarianța}

\subsection{Definiția generală}

\begin{definition}
Pentru variabile aleatoare $X$ și $Y$ cu medii finite,


\[
\mathrm{Cov}(X,Y)
= \mathbb{E}[(X - \mathbb{E}[X])(Y - \mathbb{E}[Y])].
\]


\end{definition}

\begin{proposition}[Formula dezvoltată]


\[
\mathrm{Cov}(X,Y)
= \mathbb{E}[XY] - \mathbb{E}[X]\,\mathbb{E}[Y].
\]


\end{proposition}

\subsection{Caz discret}

Dacă $(X,Y)$ ia valori $(x_i,y_j)$ cu probabilități $p_{ij}$, atunci


\[
\mathrm{Cov}(X,Y)
= \sum_{i,j} x_i y_j\, p_{ij}
- \left(\sum_i x_i p_X(x_i)\right)
  \left(\sum_j y_j p_Y(y_j)\right).
\]



\subsection{Caz continuu}

Dacă $(X,Y)$ au densitate comună $f(x,y)$, atunci


\[
\mathrm{Cov}(X,Y)
= \int_{-\infty}^{\infty}\int_{-\infty}^{\infty}
xy\, f(x,y)\, dx\, dy
- \left(\int x f_X(x)\, dx\right)
  \left(\int y f_Y(y)\, dy\right).
\]



\section{Exemple fundamentale}

\subsection{Bernoulli}

Dacă $X \sim \mathrm{Bernoulli}(p)$, atunci


\[
\mathbb{P}(X=1)=p,\quad \mathbb{P}(X=0)=1-p.
\]



\begin{proposition}
Pentru $X \sim \mathrm{Bernoulli}(p)$:


\[
\mathbb{E}[X] = p,\qquad
\mathrm{Var}(X) = p(1-p).
\]


\end{proposition}

\begin{proof}


\[
\mathbb{E}[X] = 0\cdot(1-p) + 1\cdot p = p.
\]




\[
\mathbb{E}[X^2] = 0^2(1-p) + 1^2 p = p.
\]


Deci


\[
\mathrm{Var}(X) = \mathbb{E}[X^2] - (\mathbb{E}[X])^2 = p - p^2 = p(1-p).
\]


\end{proof}

\subsection{Binomială}

Dacă $X \sim \mathrm{Bin}(n,p)$, atunci


\[
\mathbb{P}(X=k) = \binom{n}{k} p^k (1-p)^{n-k},\quad k=0,\dots,n.
\]



\begin{proposition}
Pentru $X \sim \mathrm{Bin}(n,p)$:


\[
\mathbb{E}[X] = np,\qquad
\mathrm{Var}(X) = np(1-p).
\]


\end{proposition}

\begin{proof}
Scriem $X = X_1 + \cdots + X_n$, unde $X_i \sim \mathrm{Bernoulli}(p)$ independente.
Atunci


\[
\mathbb{E}[X] = \sum_{i=1}^n \mathbb{E}[X_i] = \sum_{i=1}^n p = np.
\]


Cum $X_i$ sunt independente,


\[
\mathrm{Var}(X) = \sum_{i=1}^n \mathrm{Var}(X_i) = \sum_{i=1}^n p(1-p) = np(1-p).
\]


\end{proof}

\subsection{Geometrică}

Dacă $X \sim \mathrm{Geom}(p)$, atunci


\[
\mathbb{P}(X=k) = (1-p)^{k-1}p,\quad k=1,2,\dots
\]



\begin{proposition}
Pentru $X \sim \mathrm{Geom}(p)$:


\[
\mathbb{E}[X] = \frac{1}{p},\qquad
\mathrm{Var}(X) = \frac{1-p}{p^2}.
\]


\end{proposition}

\begin{proof}
Folosim seriile:


\[
\mathbb{E}[X] = \sum_{k=1}^\infty k(1-p)^{k-1}p = p \sum_{k=1}^\infty k q^{k-1},\ q=1-p.
\]


Se știe


\[
\sum_{k=1}^\infty k q^{k-1} = \frac{1}{(1-q)^2} = \frac{1}{p^2},
\]


deci $\mathbb{E}[X]=\frac{1}{p}$.

Pentru $\mathbb{E}[X^2]$:


\[
\mathbb{E}[X^2] = \sum_{k=1}^\infty k^2 q^{k-1} p.
\]


Se poate arăta (prin derivări succesive ale seriei geometrice) că


\[
\sum_{k=1}^\infty k^2 q^{k-1} = \frac{1+q}{(1-q)^3} = \frac{2-p}{p^3}.
\]


Deci


\[
\mathbb{E}[X^2] = p \cdot \frac{2-p}{p^3} = \frac{2-p}{p^2}.
\]


Atunci


\[
\mathrm{Var}(X) = \mathbb{E}[X^2] - (\mathbb{E}[X])^2
= \frac{2-p}{p^2} - \frac{1}{p^2} = \frac{1-p}{p^2}.
\]


\end{proof}

\subsection{Binomială negativă}

Dacă $X \sim \mathrm{NegBin}(r,p)$, atunci


\[
\mathbb{P}(X=k) = \binom{k-1}{r-1} p^r (1-p)^{k-r},\quad k=r,r+1,\dots
\]



\begin{proposition}
Pentru $X \sim \mathrm{NegBin}(r,p)$:


\[
\mathbb{E}[X] = \frac{r}{p},\qquad
\mathrm{Var}(X) = \frac{r(1-p)}{p^2}.
\]


\end{proposition}

\begin{proof}
Interpretăm $X$ ca sumă a $r$ variabile geometrice independente:


\[
X = X_1 + \cdots + X_r,\quad X_i \sim \mathrm{Geom}(p).
\]


Atunci


\[
\mathbb{E}[X] = \sum_{i=1}^r \mathbb{E}[X_i] = r \cdot \frac{1}{p} = \frac{r}{p},
\]




\[
\mathrm{Var}(X) = \sum_{i=1}^r \mathrm{Var}(X_i) = r \cdot \frac{1-p}{p^2} = \frac{r(1-p)}{p^2}.
\]


\end{proof}

\subsection{Poisson}

Dacă $X \sim \mathrm{Poisson}(\lambda)$:


\[
\mathbb{P}(X=k) = e^{-\lambda}\frac{\lambda^k}{k!},\quad k=0,1,2,\dots
\]



\begin{proposition}
Pentru $X \sim \mathrm{Poisson}(\lambda)$:


\[
\mathbb{E}[X] = \lambda,\qquad
\mathrm{Var}(X) = \lambda.
\]


\end{proposition}

\begin{proof}


\[
\mathbb{E}[X] = \sum_{k=0}^\infty k e^{-\lambda}\frac{\lambda^k}{k!}
= e^{-\lambda}\sum_{k=1}^\infty k \frac{\lambda^k}{k!}
= e^{-\lambda}\sum_{k=1}^\infty \lambda \frac{\lambda^{k-1}}{(k-1)!}
= \lambda e^{-\lambda} \sum_{j=0}^\infty \frac{\lambda^j}{j!} = \lambda.
\]


Similar, se obține $\mathbb{E}[X(X-1)] = \lambda^2$, de unde


\[
\mathrm{Var}(X) = \mathbb{E}[X^2] - (\mathbb{E}[X])^2
= \mathbb{E}[X(X-1)] + \mathbb{E}[X] - \lambda^2 = \lambda^2 + \lambda - \lambda^2 = \lambda.
\]


\end{proof}

\subsection{Hipergeometrică}

Dacă $X \sim \mathrm{Hypergeom}(N,M,n)$:


\[
\mathbb{P}(X=k) = \frac{\binom{M}{k}\binom{N-M}{n-k}}{\binom{N}{n}},\quad k=0,\dots,n.
\]



\begin{proposition}
Pentru $X \sim \mathrm{Hypergeom}(N,M,n)$:


\[
\mathbb{E}[X] = n\frac{M}{N},\qquad
\mathrm{Var}(X) = n\frac{M}{N}\left(1-\frac{M}{N}\right)\frac{N-n}{N-1}.
\]


\end{proposition}

\begin{proof}
Fie $I_i$ indicatorul că al $i$-lea element extras este succes, $i=1,\dots,n$.
Atunci $X = \sum_{i=1}^n I_i$.
Prin simetrie,


\[
\mathbb{P}(I_i=1) = \frac{M}{N},\quad \mathbb{E}[I_i]=\frac{M}{N},
\]


deci


\[
\mathbb{E}[X] = \sum_{i=1}^n \mathbb{E}[I_i] = n\frac{M}{N}.
\]


Pentru varianță, se folosește


\[
\mathrm{Var}(X) = \sum_{i=1}^n \mathrm{Var}(I_i) + 2\sum_{i<j}\mathrm{Cov}(I_i,I_j),
\]


și se obține formula cunoscută


\[
\mathrm{Var}(X) = n\frac{M}{N}\left(1-\frac{M}{N}\right)\frac{N-n}{N-1}.
\]


\end{proof}

\subsection{Multinomială}

Dacă $(X_1,\dots,X_k) \sim \mathrm{Multinomial}(n;p_1,\dots,p_k)$:


\[
\mathbb{P}(X_1=x_1,\dots,X_k=x_k)
= \frac{n!}{x_1!\cdots x_k!} p_1^{x_1}\cdots p_k^{x_k},\quad \sum_i x_i = n.
\]



\begin{proposition}
Pentru $(X_1,\dots,X_k) \sim \mathrm{Multinomial}(n;p_1,\dots,p_k)$:


\[
\mathbb{E}[X_i] = np_i,\qquad
\mathrm{Var}(X_i) = np_i(1-p_i),\qquad
\mathrm{Cov}(X_i,X_j) = -np_ip_j\ (i\neq j).
\]


\end{proposition}

\begin{proof}
Scriem $X_i = \sum_{t=1}^n I_{t,i}$, unde $I_{t,i}$ este indicatorul că a $t$-a încercare cade în categoria $i$.
Atunci


\[
\mathbb{E}[X_i] = \sum_{t=1}^n \mathbb{E}[I_{t,i}] = \sum_{t=1}^n p_i = np_i.
\]




\[
\mathrm{Var}(X_i) = \sum_{t=1}^n \mathrm{Var}(I_{t,i}) = n p_i(1-p_i).
\]


Pentru $i\neq j$,


\[
\mathrm{Cov}(X_i,X_j) = \sum_{t=1}^n \mathrm{Cov}(I_{t,i},I_{t,j})
= n\left(\mathbb{E}[I_{t,i}I_{t,j}] - p_ip_j\right).
\]


Dar $I_{t,i}I_{t,j}=0$ (o încercare nu poate fi simultan în două categorii), deci


\[
\mathrm{Cov}(X_i,X_j) = n(0 - p_ip_j) = -np_ip_j.
\]


\end{proof}

\subsection{Uniformă continuă}

Dacă $X \sim \mathrm{Unif}(a,b)$:


\[
f(x) = \frac{1}{b-a},\ x\in[a,b].
\]



\begin{proposition}
Pentru $X \sim \mathrm{Unif}(a,b)$:


\[
\mathbb{E}[X] = \frac{a+b}{2},\qquad
\mathrm{Var}(X) = \frac{(b-a)^2}{12}.
\]


\end{proposition}

\begin{proof}


\[
\mathbb{E}[X] = \int_a^b x \frac{1}{b-a}\,dx
= \frac{1}{b-a}\left.\frac{x^2}{2}\right|_a^b
= \frac{b^2-a^2}{2(b-a)} = \frac{a+b}{2}.
\]




\[
\mathbb{E}[X^2] = \int_a^b x^2 \frac{1}{b-a}\,dx
= \frac{1}{b-a}\left.\frac{x^3}{3}\right|_a^b
= \frac{b^3-a^3}{3(b-a)} = \frac{a^2+ab+b^2}{3}.
\]


Deci


\[
\mathrm{Var}(X) = \mathbb{E}[X^2] - (\mathbb{E}[X])^2
= \frac{a^2+ab+b^2}{3} - \frac{(a+b)^2}{4}
= \frac{(b-a)^2}{12}.
\]


\end{proof}

\subsection{Exponențială}

Dacă $X \sim \mathrm{Exp}(\lambda)$:


\[
f(x) = \lambda e^{-\lambda x},\ x\ge0.
\]



\begin{proposition}
Pentru $X \sim \mathrm{Exp}(\lambda)$:


\[
\mathbb{E}[X] = \frac{1}{\lambda},\qquad
\mathrm{Var}(X) = \frac{1}{\lambda^2}.
\]


\end{proposition}

\begin{proof}


\[
\mathbb{E}[X] = \int_0^\infty x \lambda e^{-\lambda x}\,dx.
\]


Cu $t=\lambda x$,


\[
\mathbb{E}[X] = \frac{1}{\lambda}\int_0^\infty t e^{-t}\,dt = \frac{1}{\lambda}\Gamma(2) = \frac{1}{\lambda}.
\]


Similar,


\[
\mathbb{E}[X^2] = \int_0^\infty x^2 \lambda e^{-\lambda x}\,dx
= \frac{1}{\lambda^2}\int_0^\infty t^2 e^{-t}\,dt = \frac{1}{\lambda^2}\Gamma(3) = \frac{2}{\lambda^2}.
\]


Deci


\[
\mathrm{Var}(X) = \mathbb{E}[X^2] - (\mathbb{E}[X])^2
= \frac{2}{\lambda^2} - \frac{1}{\lambda^2} = \frac{1}{\lambda^2}.
\]


\end{proof}

\subsection{Gamma}

Dacă $X \sim \mathrm{Gamma}(\alpha,\lambda)$:


\[
f(x) = \frac{\lambda^\alpha}{\Gamma(\alpha)} x^{\alpha-1} e^{-\lambda x},\ x\ge0.
\]



\begin{proposition}
Pentru $X \sim \mathrm{Gamma}(\alpha,\lambda)$:


\[
\mathbb{E}[X] = \frac{\alpha}{\lambda},\qquad
\mathrm{Var}(X) = \frac{\alpha}{\lambda^2}.
\]


\end{proposition}

\begin{proof}


\[
\mathbb{E}[X] = \int_0^\infty x \frac{\lambda^\alpha}{\Gamma(\alpha)} x^{\alpha-1} e^{-\lambda x}\,dx
= \frac{\lambda^\alpha}{\Gamma(\alpha)}\int_0^\infty x^\alpha e^{-\lambda x}\,dx.
\]


Cu $t=\lambda x$,


\[
\mathbb{E}[X] = \frac{\lambda^\alpha}{\Gamma(\alpha)}\frac{1}{\lambda^{\alpha+1}}\int_0^\infty t^\alpha e^{-t}\,dt
= \frac{1}{\lambda}\frac{\Gamma(\alpha+1)}{\Gamma(\alpha)} = \frac{\alpha}{\lambda}.
\]


Analog,


\[
\mathbb{E}[X^2] = \frac{\lambda^\alpha}{\Gamma(\alpha)}\int_0^\infty x^{\alpha+1} e^{-\lambda x}\,dx
= \frac{1}{\lambda^2}\frac{\Gamma(\alpha+2)}{\Gamma(\alpha)} = \frac{\alpha(\alpha+1)}{\lambda^2}.
\]


Deci


\[
\mathrm{Var}(X) = \mathbb{E}[X^2] - (\mathbb{E}[X])^2
= \frac{\alpha(\alpha+1)}{\lambda^2} - \frac{\alpha^2}{\lambda^2}
= \frac{\alpha}{\lambda^2}.
\]


\end{proof}

\subsection{Normală}

Dacă $X \sim \mathcal{N}(\mu,\sigma^2)$:


\[
f(x) = \frac{1}{\sigma\sqrt{2\pi}}\exp\left(-\frac{(x-\mu)^2}{2\sigma^2}\right).
\]



\begin{proposition}
Pentru $X \sim \mathcal{N}(\mu,\sigma^2)$:


\[
\mathbb{E}[X] = \mu,\qquad
\mathrm{Var}(X) = \sigma^2.
\]


\end{proposition}

\begin{proof}
Scriem $X = \mu + \sigma Z$, unde $Z \sim \mathcal{N}(0,1)$.
Atunci


\[
\mathbb{E}[X] = \mu + \sigma \mathbb{E}[Z].
\]


Densitatea lui $Z$ este pară, deci $z\varphi(z)$ este impară și


\[
\mathbb{E}[Z] = \int_{-\infty}^\infty z\varphi(z)\,dz = 0,
\]


deci $\mathbb{E}[X]=\mu$.

Pentru varianță,


\[
\mathrm{Var}(X) = \mathrm{Var}(\mu + \sigma Z) = \sigma^2 \mathrm{Var}(Z).
\]


Se știe că $\mathrm{Var}(Z)=1$, deci $\mathrm{Var}(X)=\sigma^2$.
\end{proof}


\end{document}
